% Page-breaking rules for the PDF edition.
%
% LaTeX's defaults happily strand a paragraph's last line at the top of an
% otherwise empty page, or leave a heading alone at the foot of one. These
% penalties forbid both, and \raggedbottom lets a page simply end short
% instead of stretching the vertical space to justify a ragged bottom.

% Never split a single line of a paragraph away from the rest, in either
% direction: no lone last line on a new page, no lone first line at the foot.
\widowpenalty=10000
\clubpenalty=10000
\displaywidowpenalty=10000

% Never leave a heading at the bottom of a page with its text overleaf.
% The starred option also moves headings that land just above the bottom margin.
\usepackage[nobottomtitles*]{titlesec}

% Prefer a page that ends short over one padded with stretched spacing.
\raggedbottom

% Widow and orphan penalties do not govern bulleted lists, which is how a lone
% final bullet ends up stranded at the top of an otherwise blank page. These
% are LaTeX internals, hence \makeatletter. The values discourage a break
% before, inside, and after a list strongly, without forbidding it outright:
% a genuinely long list must still be allowed to span two pages.
\makeatletter
\@beginparpenalty=9000
\@itempenalty=9000
\@endparpenalty=9000
\makeatother

% This book is full of bulleted lists. Tightening them slightly is both better
% typography for a reference text and fewer lists that overflow a page by a
% single item, which is the case penalties alone cannot solve on a full page.
\usepackage{enumitem}
\setlist[itemize]{itemsep=2pt, parsep=0pt, topsep=4pt}
\setlist[enumerate]{itemsep=2pt, parsep=0pt, topsep=4pt}
